\section{Verification purpose}

\textbf{Case ID:} \texttt{\CaseID}. This case verifies that ELMFIRE's raster indexing and level-set transport are covariant under a \RotationAngle{} rotation. It is the portable successor to Adam Laird's \texttt{inputs-orientation} scenario.

\section{Mathematical formulation and numerical implementation}

For homogeneous isotropic forcing, the operator $\mathcal{L}$ in $\partial_t\phi=\mathcal{L}(\phi)$ commutes with any rigid rotation $Q$: $\mathcal{L}(\phi\circ Q^{-1})=(\mathcal{L}\phi)\circ Q^{-1}$. A square north-up raster rotated by exactly 90 degrees requires no interpolation, so cell-by-cell comparison is available.

\section{Derivation of expected behavior}

The rotated variant is generated with \texttt{numpy.rot90} from the asymmetric L-shaped initial negative-$\phi$ region. Fuel, moisture, slope, and wind are spatially uniform. Therefore the rotated arrival-time validity mask and values must equal a 90-degree rotation of the original result, apart from bounded grid-level roundoff.

% BEGIN EXPLICIT METRIC DERIVATION
The reference array is generated directly from the original ELMFIRE
arrival-time result by a 90-degree counterclockwise rotation,
\(T_{\mathrm{ref}}=\operatorname{rot90}(T_{\mathrm{original}})\); no
arrival values are invented.  With
\(V=\{i:T_{\mathrm{ref},i}\text{ and }T_{\mathrm{rot},i}\text{ finite}\}\)
and \(U=\{i:\text{either value is finite}\}\), finite-mask agreement is
\[
 A_m=\frac{|V|}{|U|},
\]
whose exact rotational reference is 1 and whose substituted test is
\(A_m\geq0.99\).  On the mutual finite mask, the value metric is
\[
 E_T=\frac{|V|^{-1}\sum_{i\in V}
 |T_{\mathrm{ref},i}-T_{\mathrm{rot},i}|}
 {\max[P_{90}(T_{\mathrm{ref},i}:i\in V),1~\mathrm{s}]},
\]
with exact reference zero and acceptance substitution \(E_T\leq0.02\).
The asymmetric generated ignition mask makes this comparison sensitive to
an erroneous reflection or transposition.
% END EXPLICIT METRIC DERIVATION

\section{Verification metrics and acceptance criteria}

Finite-mask intersection over union must be at least 0.99. On mutually finite cells, normalized mean absolute arrival-time error must be at most 0.02. Mask agreement detects row/column reversal and georeferencing errors; value error detects orientation-dependent numerical evolution.

\section{Simulation configuration and predicted outcome}

Both variants use a 2000 m square domain at \CellSize{}, Fuel Model 10, flat terrain, zero wind, and identical timing. The only changed scientific input is the exact raster rotation of \texttt{phi.tif}. The second result is predicted to be the corresponding rotation of the first.

\begin{center}\begin{tabular}{p{0.22\textwidth}p{0.22\textwidth}p{0.12\textwidth}p{0.34\textwidth}}\toprule
Parameter & Configured value & Units & Role \\
\midrule Domain/grid & 2000 m square; 5 & m/cell & Exact square raster rotation \\
Fuel/forcing & FM10; flat; calm & --- & Isotropic homogeneous operator \\
Time control & step 1, maximum 5; stop 7200 & s & Common exact comparison time \\
Initial field & asymmetric L shape & --- & Makes orientation errors observable \\
Variant transform & 90-degree rotation & degrees & Exact no-interpolation covariance test \\
\bottomrule\end{tabular}\end{center}

Figure~\ref{fig:input-configuration} is generated mechanically from the prepared case-local GeoTIFF inputs. It identifies the selected source artifact and representative variant; the left panel shows a fuel, terrain, moisture, wind, canopy, or other controlling field selected for this experiment, and the right panel shows the cells selected by the non-positive initial level-set field. No values in this figure are redrawn or inferred from the expected response.

\begin{figure}[H]
  \centering
  \IfFileExists{../figures/input_configuration.pdf}{%
    \includegraphics[width=0.98\textwidth,height=0.46\textheight,keepaspectratio]{input_configuration.pdf}%
  }{\fbox{\parbox{0.88\textwidth}{The prepared-input figure has not been generated. Run the case-local preprocessing and spatial-evidence step.}}}
  \caption{Prepared input conditions from case-local generated rasters. Coordinates are in metres; the colorbar identifies the displayed field and its physical units or category. The ignition panel is derived directly from \(\phi\leq0\).}
  \label{fig:input-configuration}
\end{figure}

\section{Actual simulation results}

The generated status is \textbf{\VerificationStatus}; \PassingMetricCount{} of \MetricCount{} metrics pass across \VariantCount{} variants. Missing output produces NOT EVALUABLE.
\begin{figure}[H]\centering\IfFileExists{../figures/verification_summary.pdf}{\includegraphics[width=\textwidth]{verification_summary.pdf}}{\fbox{No result figure is available.}}\caption{Original, independently rotated, and difference arrival-time fields.}\end{figure}

\section{Calculated metrics and verification decision}

\begin{center}\begin{tabular}{@{}llll@{}}\toprule Metric & Expected & Calculated & Status \\\midrule \MetricRows \bottomrule\end{tabular}\end{center}
The authoritative decision is \textbf{\VerificationStatus} from \texttt{outputs/metrics.json}.

\section{Technical interpretation}

Mask disagreement indicates inconsistent raster orientation, bounds, or row/column indexing. Value disagreement with matching masks indicates a directional bias in transport or output serialization.

\section{Scope and limitations}

The 90-degree test avoids resampling and therefore does not assess arbitrary-angle reprojection. It verifies executable behavior, not the internal bilinear-interpolation routine in isolation.
